
\begin{proposition}[Hopf-axis adjacency with arbitrary multiplicities and point blowups]\label{cdp:st:C.2.UN.40}
Let $f:U\to\Delta$ be a proper degeneration of smooth complex
two-tori with smooth total space, and suppose that its reduced
central support has simple normal crossings. Write its full
fibre as $\sum_i m_iE_i$. Choose a generator of the rank-one
sheaf $f_*K_U$ and write its divisor as $\sum_i k_iE_i$.
Set $r_i=(k_i+1)/m_i$.

Suppose that the minimal model $H_i$ of a component $E_i$
is a Hopf surface of algebraic dimension zero. The contraction
$\rho_i:E_i\to H_i$ may consist of arbitrarily many point
blowdowns. Choose an actual finite unramified primary Hopf
cover $p:P\to H_i$.

If $P$ has diagonal contraction $(z,w)\mapsto(\alpha z,\beta w)$,
the images of both coordinate axes are boundary curves on
$H_i$: their strict transforms in $E_i$ meet another central
component. For each such strict transform, let $E_j$ be its
unique adjacent component at the generic point. Then
\[
                 m_i(k_j+1)-m_j(k_i+1)=0,
                 \qquad r_j=r_i.
\tag{C.HA.1}\label{cdp:eq:C.HA.1}
\]
If $P$ has non-diagonal contraction
$(z,w)\mapsto(\beta^r z+w^r,\beta w)$, $r\geq1$, the image
of its unique elliptic curve is a boundary curve. Its unique
generic neighbour satisfies
\[
       m_i(k_j+1)-m_j(k_i+1)=-m_i r,
       \qquad m_j(r_j-r_i)=-r<0.
\tag{C.HA.2}\label{cdp:eq:C.HA.2}
\]
These are integral divisor-class consequences on the actual
components. In particular a non-diagonal algebraic-dimension-zero
Hopf component cannot meet another algebraic-dimension-zero
Hopf component along this elliptic curve. Diagonal such
components can meet along their axes only at equal ratios.
\end{proposition}

\begin{proof}
The direct image $f_*K_U$ is coherent, torsion-free and of
rank one on the disc, because a smooth torus fibre has one
holomorphic two-form. Shrink the disc and choose a generator.
Its evaluation is an actual holomorphic three-form on $U$.
On every punctured fibre it is a nonzero translation volume
form, so its zero divisor is vertical and gives
$K_U\simeq\mathcal O_U(\sum_i k_iE_i)$, with $k_i\geq0$.
For $j\ne i$ put $D_{ij}=E_j|_{E_i}$ and
$D_i=\sum_{j\ne i}D_{ij}$. Each is the reduced intersection
divisor specified by the simple-normal-crossing coordinates.
The full fibre is principal. Adjunction and its restriction
to $E_i$ therefore give the following integral identity:
\[
 m_i(K_{E_i}+D_i)
 \sim\sum_{j\ne i}
       \bigl[m_i(k_j+1)-m_j(k_i+1)\bigr]D_{ij}.
\tag{C.HA.3}\label{cdp:eq:C.HA.3}
\]
Indeed, first use
$K_{E_i}=((k_i+1)E_i+\sum_{j\ne i}k_jE_j)|_{E_i}$,
then replace $m_i E_i|_{E_i}$ by
$-\sum_{j\ne i}m_jD_{ij}$, and finally add $m_iD_i$.
Thus \eqref{cdp:eq:C.HA.3} has no rational divisor or divisibility convention.

We specify its pushforward through the point blowdowns.
For a proper birational map of smooth surfaces the common
meromorphic function field sends the divisor of a function
to the divisor of the same function. Sending every prime
curve dominating a curve to that curve and every exceptional
prime curve to zero consequently induces a homomorphism
of divisor-class groups. It sends $K_{E_i}$ to $K_{H_i}$:
pull back a meromorphic two-form on $H_i$ and observe that
its additional Jacobian divisor is exceptional. This is a
divisor-class homomorphism; no equality with the coherent
direct image of an arbitrary line bundle is asserted.
Applying it to \eqref{cdp:eq:C.HA.3} gives
\[
 m_i(K_{H_i}+B_i)
 \sim\sum_{j\ne i}
       \bigl[m_i(k_j+1)-m_j(k_i+1)\bigr]\rho_{i*}D_{ij},
 \qquad B_i=\rho_{i*}D_i.
\tag{C.HA.4}\label{cdp:eq:C.HA.4}
\]
The divisor $B_i$ is reduced. A point-blowup morphism is an
isomorphism at the generic point of every curve on its target,
so there is exactly one strict transform dominating that curve.
At its generic point the SNC equation has exactly two branches,
$x_i x_j=0$; an additional branch can meet it only at isolated
triple points. Hence this strict transform occurs in exactly
one $D_{ij}$, with coefficient one.

The pullback of \eqref{cdp:eq:C.HA.4} through the specified unramified
cover $p:P\to H_i$ is an identity of integral divisor classes,
and $p^*K_{H_i}=K_P$. We record precisely the curve and line
classes on $P$. The primary normal forms are those in
\cite[Theorem~2.1]{Khaled2024Hopf}. The finite meromorphic
field comparison proves $a(P)=a(H_i)=0$, by taking elementary
symmetric functions on the finite generic sheets, as in
(C.CF.9).

Here is also a direct check of the curves needed in that
comparison. The inverse image of a compact curve in
$\mathbb C^2\setminus\{0\}$ has analytic closure across
the origin by the Remmert--Stein extension theorem. Its
divisor on $\mathbb C^2$ is principal, say $(F=0)$.
For the contraction $\gamma$ one has
$F\circ\gamma=uF$ with $u$ a holomorphic unit. Write
$u=u(0)e^h$, $h(0)=0$. The locally uniformly convergent sum
$a=\sum_{n\geq0}h\circ\gamma^n$ satisfies
$a-a\circ\gamma=h$. Replacing $F$ by $e^aF$ makes it an
eigenfunction of $\gamma^*$ with constant eigenvalue.

For a diagonal contraction, distinct monomials can have
the same eigenvalue only when $\alpha^a\beta^b=1$ for
some nonzero integer pair. Because $|\alpha|,|\beta|<1$,
such a pair has opposite signs and gives a nonconstant
invariant meromorphic monomial, contrary to $a(P)=0$.
Thus each eigenfunction is one monomial, and the only
curves are the two axes $C_1,C_2$. In the non-diagonal
case use weighted degree $\deg z=r$, $\deg w=1$.
On the finite-dimensional space of weighted degree $d$
the contraction has eigenvalue $\beta^d$ times the
translation operator on polynomials in $z/w^r$.
Its eigenvectors are precisely the constants in that
variable. The eigenvalues $\beta^d$ for different $d$
are distinct, so every eigenfunction is a power of $w$.
There is exactly the curve $C=(w=0)$.

For completeness, denote by $L_\lambda$ the line bundle
whose lifted fibre coordinate is multiplied by $\lambda$
under the contraction. The map
$\lambda\mapsto L_\lambda$ is injective. A trivialization
would be a nowhere-zero holomorphic eigenfunction on
$\mathbb C^2\setminus\{0\}$. Hartogs extends it and its
inverse across zero, and evaluation there forces
$\lambda=1$. Coordinate functions and the Jacobian
therefore give the exact classes
\[
\begin{array}{ll}
 \mathcal O(C_1)=L_\alpha,
 \quad\mathcal O(C_2)=L_\beta,
 &K_P=\mathcal O(-C_1-C_2),\\[2pt]
 \mathcal O(C)=L_\beta,
 &K_P=\mathcal O(-(r+1)C)
 \quad\hbox{in the non-diagonal case}.
\end{array}
\tag{C.HA.5}\label{cdp:eq:C.HA.5}
\]
The first two classes are integrally independent: a relation,
or even a torsion relation, would give the forbidden
positive-power resonance after raising to a power.
The class of $C$ in the second line has infinite order
because $0<|\beta|<1$.

Every curve on $H_i$ has nonempty reduced inverse image
on $P$, since $p$ is unramified and surjective. Thus the
pullback of $B_i$ is a reduced subset of the indicated
axis curves. In the diagonal case, compare coefficients
of their independent classes in the pullback of \eqref{cdp:eq:C.HA.4}.
If an axis were absent, its coefficient on the left
would be $-m_i$ and on the right zero. Each axis is
therefore present. Its coefficient on the left is then
zero, and the coefficient on the right is exactly that
of its unique generic neighbour. This proves \eqref{cdp:eq:C.HA.1},
including when one boundary curve has more than one
component in its finite inverse image.

In the non-diagonal case, absence of $C$ would instead
give the impossible nonzero class $-m_i(r+1)C\sim0$.
It is present, with coefficient one. Comparing its
infinite-order class gives
\[
       -m_i r=m_i(k_j+1)-m_j(k_i+1),
\tag{C.HA.6}\label{cdp:eq:C.HA.6}
\]
which is \eqref{cdp:eq:C.HA.2}. This proves both assertions on the
original component after all its point blowups, with
the actual finite-primary-cover comparison included.

Finally, the shared strict axis is a smooth elliptic
curve. It cannot be contracted by a point-blowdown
morphism on the neighbouring smooth component, since
all exceptional curves of such a morphism are rational.
If that neighbour also has an algebraic-dimension-zero
Hopf minimal model, this same curve is one of its
minimal axes. Its diagonal rule would force equal
ratios, contradicting \eqref{cdp:eq:C.HA.2}; its non-diagonal rule
would force the opposite strict inequality. Neither
is possible. The diagonal--diagonal case merely repeats
the equality \eqref{cdp:eq:C.HA.1}. The theorem gives these exact
adjacency restrictions and does not infer a global
nonexistence assertion from a local decrease.
\end{proof}


